\documentclass[hidelinks]{article}
\usepackage{stex}

\libinput{preambleCS}

\title{Deterministic and Nondeterministic Finite State Automata}
\author{}
\date{}

\begin{document}

\maketitle

\begin{smodule}[title={Deterministic and Nondeterministic Finite State Automata}]{DeterministicAndNondeterministicFiniteAutomata}
\usemodule{ComputerScience/FormalLanguagesAndAutomata?IntroductionToFormalLanguagesAndAutomata}

\symdecl*{deterministic finite-state automaton}
\symdecl*{extended transition function}
\symdecl*{nondeterministic finite-state automaton}
\symdecl*{transition relation}
\symdecl*{lambda-transition}
\symdecl*{junk state}

\section{Overview}

The previous lecture introduced the \sn{finite-state automaton} as the simplest of our automaton models, with no auxiliary storage. This note formalises the FSA in two flavours: the \sn{deterministic finite-state automaton} (DFA), in which exactly one transition is defined at each (state, input) pair, and the \sn{nondeterministic finite-state automaton} (NDFA), in which a state may have several transitions on the same input symbol -- or none at all -- and may also use \sns{lambda-transition} that consume no input. Both flavours recognise exactly the same class of languages, the \sns{regular language}, but the nondeterministic model is often more compact and natural to write down.

\section{Deterministic finite state automata}

\begin{sdefinition}[for=deterministic finite-state automaton]
A \definame{deterministic finite-state automaton} (DFA) is a 5-tuple $M = \langle Q, \Sigma, \delta, i, F \rangle$, where
\begin{itemize}
  \item $Q$ is a finite set of \emph{states};
  \item $\Sigma$ is the \sn{alphabet}, a finite set of possible input symbols;
  \item $\delta : Q \times \Sigma \to Q$ is the \emph{transition function} (a total function);
  \item $i \in Q$ is the \emph{initial state};
  \item $F \subseteq Q$ is the set of \emph{final} (accepting) states.
\end{itemize}
\end{sdefinition}

\subsection{Extended transition function}

To talk about acceptance of whole strings rather than single symbols, we extend $\delta$ from inputs in $\Sigma$ to inputs in $\Sigma^*$.

\begin{sdefinition}[for=extended transition function]
The \definame{extended transition function} $\overline{\delta} : Q \times \Sigma^* \to Q$ of a \sn{deterministic finite-state automaton} is recursively defined by
\[
\overline{\delta}(q, \lambda) = q, \qquad \overline{\delta}(q, aw) = \overline{\delta}(\delta(q, a), w)
\]
for $a \in \Sigma$ and $w \in \Sigma^*$. Some textbooks write $\widehat{\delta}$ or $\delta^*$ in place of $\overline{\delta}$.
\end{sdefinition}

\subsection{The language accepted by a DFA}

A DFA $M = \langle Q, \Sigma, \delta, i, F \rangle$ \emph{accepts} a \sn{string} $w$ if $\overline{\delta}(i, w) \in F$. Note that the empty string $\lambda$ is accepted by $M$ if and only if $i \in F$. The \emph{language accepted by $M$} is the set of all strings that $M$ accepts:
\[
L(M) = \{ w \in \Sigma^* \mid \overline{\delta}(i, w) \in F \}.
\]

\begin{example}[A DFA for ``equal counts of $0$s and $1$s, both even'']
We construct a DFA $M$ with $\Sigma = \{0, 1\}$ such that
\[
L(M) = \{ w \in \Sigma^* \mid \text{the numbers of $0$s and $1$s in $w$ are both even}\}.
\]
The idea is that $M$ tracks, with each symbol it reads, whether the numbers of $0$s and $1$s so far are even or odd. This suggests a four-state machine with $q_0$ (even $0$s, even $1$s), $q_1$ (odd $0$s, even $1$s), $q_2$ (even $0$s, odd $1$s), $q_3$ (odd $0$s, odd $1$s); $q_0$ is the initial and only accepting state. Each input symbol toggles the corresponding parity:

\begin{automaton}{2.6cm}
  \startfinalstate (q0)              {$q_0$} ;
  \state          (q1) [right of=q0] {$q_1$} ;
  \state          (q2) [below of=q0] {$q_2$} ;
  \state          (q3) [right of=q2] {$q_3$} ;
  \path (q0) edge [bend left]  node {$0$} (q1)
        (q1) edge [bend left]  node {$0$} (q0)
        (q0) edge [bend left]  node [left] {$1$} (q2)
        (q2) edge [bend left]  node [right] {$1$} (q0)
        (q1) edge [bend left]  node [left] {$1$} (q3)
        (q3) edge [bend left]  node [right] {$1$} (q1)
        (q2) edge [bend left]  node [below] {$0$} (q3)
        (q3) edge [bend left]  node [above] {$0$} (q2) ;
\end{automaton}
\end{example}

\subsection{Drawing conventions}

Two graphical conventions reduce diagram clutter without changing the underlying machine.

\paragraph{Junk states.} Each DFA state must have an outgoing transition for every symbol in $\Sigma$, because $\delta$ is a total function. This often forces a \emph{junk} (or sink) state with no path to a final state.

\begin{sdefinition}[for=junk state]
A \definame{junk state} (or \emph{sink state}) of a \sn{deterministic finite-state automaton} is a non-accepting state $q$ from which no final state is reachable, so that any computation entering $q$ cannot subsequently accept.
\end{sdefinition}

For simplicity, we agree that \sns{junk state} and their incoming transitions can be omitted in transition diagrams; the missing transitions are implicitly understood to lead to a junk state.

\paragraph{Stacked transition labels.} If two or more transitions exist between the same pair of states, the labels can be stacked on a single arrow rather than drawn as parallel arrows.

\section{Nondeterministic finite state automata}

By definition, every state of a DFA has a unique successor for every input symbol. \sns{nondeterministic finite-state automaton} (NDFAs) relax this restriction in two ways:
\begin{itemize}
  \item a state may have \emph{several} outgoing transitions with the same label, and the machine \emph{chooses} which one to execute; and
  \item a state may have \emph{no} outgoing transition for some input symbol -- if the machine reads that symbol, it simply stops, so NDFAs do not need junk states.
\end{itemize}
Formally, the deterministic transition function $\delta$ is replaced by a relation, the \sn{transition relation}.

\subsection{Lambda transitions}

NDFAs may additionally contain \sns{lambda-transition}, which consume no input symbol.

\begin{sdefinition}[for=lambda-transition]
A \definame{lambda-transition} (also called a $\lambda$-transition or $\varepsilon$-transition) of a \sn{nondeterministic finite-state automaton} is a transition that changes state without consuming an input symbol; in transition diagrams it is drawn with the label $\lambda$.
\end{sdefinition}

\subsection{Formal definition of an NDFA}

\begin{sdefinition}[for=nondeterministic finite-state automaton]
A \definame{nondeterministic finite-state automaton} (NDFA) is a 5-tuple $M = \langle Q, \Sigma, \rho, i, F \rangle$, where $Q$, $\Sigma$, $i$, $F$ are as for a \sn{deterministic finite-state automaton} and
\[
\rho \subseteq 2^{Q \times (\Sigma \cup \{\lambda\}) \times Q}
\]
is the \emph{transition relation}.
\end{sdefinition}

\begin{sdefinition}[for=transition relation]
The \definame{transition relation} $\rho$ of an NDFA contains a triple $(p, a, q)$ exactly when, in state $p$ reading $a$ (or $a = \lambda$, taking a $\lambda$-transition), the machine may move to state $q$. Some textbooks instead use a set-valued transition function $\delta_{NDFA} : Q \times (\Sigma \cup \{\lambda\}) \to 2^Q$, which conveys the same information.
\end{sdefinition}

\subsection{Extended transition relation}

Extending $\rho$ to strings is more involved than for DFAs because of the $\lambda$-transitions. We first define a relation $\rho^* \subseteq 2^{Q \times \Sigma \times Q}$ that captures ``state $q$ reachable from $p$ by reading $a$ amid (possibly empty) sequences of $\lambda$-transitions before and after the $a$-transition'':
\begin{enumerate}
  \item if $(p, a, q) \in \rho$ then $(p, a, q) \in \rho^*$;
  \item if $(p, a, r) \in \rho^*$ and $(r, \lambda, q) \in \rho$ then $(p, a, q) \in \rho^*$;
  \item if $(p, \lambda, r) \in \rho$ and $(r, a, q) \in \rho^*$ then $(p, a, q) \in \rho^*$.
\end{enumerate}
The \emph{extended transition relation} $\overline{\rho} \subseteq 2^{Q \times \Sigma^* \times Q}$ is then defined by
\begin{enumerate}
  \item $(q, \lambda, r) \in \overline{\rho}$ if $r$ is reachable from $q$ by a (possibly empty) sequence of $\lambda$-transitions;
  \item $(q, aw, r) \in \overline{\rho}$ if there exists a state $p$ such that $(q, a, p) \in \rho^*$ and $(p, w, r) \in \overline{\rho}$.
\end{enumerate}

\subsection{The language accepted by an NDFA}

An NDFA $M = \langle Q, \Sigma, \rho, i, F \rangle$ \emph{accepts} a string $w \in \Sigma^*$ if there exists a state $q \in F$ such that $(i, w, q) \in \overline{\rho}$. The empty string is accepted by $M$ iff some final state is reachable from $i$ by a (possibly empty) sequence of $\lambda$-transitions. The \emph{language accepted by $M$} is
\[
L(M) = \{ w \in \Sigma^* \mid (i, w, q) \in \overline{\rho} \text{ for some } q \in F \}.
\]

\begin{example}[A small NDFA over $\{a, b\}$]
The following NDFA accepts the language $\{ a a^n a \mid n \geq 0 \} \cup \{ a a^n b a \mid n \geq 0 \} \cup \cdots$ -- informally, it reads an $a$, loops on $a$ or $b$ in $q_1$, and then reads a final $a$:
\begin{automaton}{3.2cm}
  \startstate (q0) {$q_0$} ;
  \state      (q1) [right of=q0] {$q_1$} ;
  \finalstate (q2) [right of=q1] {$q_2$} ;
  \path (q0) edge              node {$a$} (q1)
        (q1) edge [loop above] node [align=center] {$a$\\$b$} (q1)
        (q1) edge              node {$a$} (q2) ;
\end{automaton}
The accepted language is $\{ a w a \mid w \in \{a, b\}^* \}$, expressed compactly thanks to nondeterminism: at $q_1$, the machine can either loop on the next symbol or commit to the final $a$ transition.
\end{example}

\subsection{Why allow nondeterminism?}

Digital computers are deterministic, so why allow nondeterminism in our model? Three reasons:
\begin{itemize}
  \item Some languages can be specified more compactly using nondeterministic automata, as the example above shows.
  \item When we apply algorithms to create or modify FSAs (for instance, the constructions for closure properties in a later note), the results are often nondeterministic.
  \item Even though NDFAs look more powerful, the next note shows that they accept exactly the same class of languages as DFAs -- the \sns{regular language}.
\end{itemize}

\section{Summary}

\sns{deterministic finite-state automaton} have a total transition function $\delta$ that is extended to strings by the \sn{extended transition function} $\overline{\delta}$, accepting a string when the run ends in a final state. \sns{nondeterministic finite-state automaton} replace $\delta$ with a \sn{transition relation} $\rho$ that admits multiple successors, missing transitions, and \sns{lambda-transition}, accepting a string when at least one branch of computation finishes in a final state. The next note shows that NDFAs and DFAs are equivalent: every NDFA can be turned into a DFA accepting the same language by the subset construction.

\end{smodule}

\end{document}
